\chapter{Behavioral Primes}
\label{sec:combinators}
% ===================================================================

Chapter~\ref{ch:igslt} gave the definition of an interactive GSLT --- a
theory with a designated site $\inter$ at which two terms meet, a base
rule that fires there without contextual hypothesis, and context rules
saying when the base rule may propagate --- and argued that the symmetry
or asymmetry of that site is a real ontological distinction rather than a
notational one.
The argument of this part needs one further thing from that structure,
and this short chapter supplies it.

The question is whether an interactive theory has \emph{atoms}: a
supply of terms out of which everything else can be built by interaction
alone.
If it does, then an agent doing science on such a theory has something to
find, and the population arguments of the chapters that follow have a
target.
If it does not, there is no bottom for the scientific method to converge
to, and the whole question of what an agent can and cannot discover
changes character.

\section{Behavioral Primes}

\begin{definition}[Combinator Set]
A \emph{combinator set} for an interactive GSLT $\GSLT$ is a set
$\comb \subseteq \terms(\GSLT)$ such that every term
$P \in \terms(\GSLT)$ is bisimilar to a finite interaction of
elements of $\comb$:
\[
  \forall P \in \terms(\GSLT),\;
  \exists c_1, \ldots, c_n \in \comb \text{ such that }
  P \bisim c_1 \inter c_2 \inter \cdots \inter c_n.
\]
\end{definition}

\begin{definition}[Behavioral Prime]
An element $c \in \comb$ is a \emph{behavioral prime} if it is not
bisimilar to any interaction $c' \inter c''$ of smaller terms.
A combinator set $\comb$ is \emph{prime} if every element is a
behavioral prime and no proper subset of $\comb$ is still a
combinator set.
\end{definition}

\begin{remark}[Computational Quarks]
The elements of a prime combinator set $\comb$ are the
\emph{computational quarks} of $\GSLT$: irreducible behavioral units
whose interactions generate all observable behavior.
They are \emph{not} the bisimulation equivalence classes, which may
be equivalence classes of arbitrarily complex composite behavior.
The primes are indivisible in the interaction-factorization sense:
they cannot be split into independently interacting parts.
\end{remark}

\begin{conjecture}[Existence of Prime Combinator Sets]
Every interactive GSLT with a finitely generated grammar admits a
prime combinator set.
\end{conjecture}

\begin{remark}
Unique factorization --- the analogue of the fundamental theorem of
arithmetic for behavioral primes --- is a stronger condition that
need not hold in general.
Its failure in a given GSLT would be physically meaningful:
it would correspond to composite objects that cannot be assigned a
unique quark content, analogous to the situation in QCD where
the quark content of hadrons is not always uniquely determined.
Whether unique factorization holds is a structural question about
the specific GSLT.
\end{remark}

% ===================================================================
